% concatenated from main.tex
\documentclass{amsart}

\usepackage{xcolor}


\setcounter{secnumdepth}{5}

\usepackage{amsfonts}
\usepackage{amssymb}
\usepackage{graphicx}
\usepackage{mathrsfs}
\usepackage{amscd}
\usepackage{hyperref}
\usepackage{stmaryrd}
\usepackage{caption}
\usepackage{subcaption}
\usepackage{color}
\usepackage{xcolor}
\usepackage{comment}
\usepackage{bm,bbm}
\usepackage{stmaryrd}
\usepackage{euscript}
\usepackage{microtype}
\usepackage{enumitem}

%Newly-Added packages
\usepackage{tikz} 
\usepackage{tikz-cd}
\usepackage{mathtools}

\newtheorem{theorem}{Theorem}[section]
\newtheorem{lemma}[theorem]{Lemma}
\newtheorem{corollary}[theorem]{Corollary}
\newtheorem{setup}[theorem]{Setup}
\newtheorem{prop}[theorem]{Proposition}
\newtheorem{defn}[theorem]{Definition}
\newtheorem{example}[theorem]{Example}
\newtheorem{remark}[theorem]{Remark}
\newtheorem{conjecture}[theorem]{Conjecture}
\newtheorem{assumption}[theorem]{Assumption}
\newtheorem{problem}[theorem]{Problem}
\newtheorem{claim}[theorem]{Claim}
\newtheorem{construction}[theorem]{Construction}

\numberwithin{equation}{section}




\newcommand{\cat}{
	\mathscr
}
\newcommand{\N}{\mathbb{N}}
\newcommand{\Z}{\mathbb{Z}}
\newcommand{\R}{\mathbb{R}}
\newcommand{\U}{\mathbb{U}}
\newcommand{\C}{\mathbb{C}}
\newcommand{\bS}{\mathbb{S}}
\newcommand{\bi}{\mathbf{i}}
\newcommand{\bT}{\mathbf{T}}
\newcommand{\bb}{\boldsymbol{b}}
\newcommand{\bv}{\boldsymbol{v}}
\newcommand{\bx}{\boldsymbol{x}}
\newcommand{\WH}[1]{\widehat{#1}}
\newcommand{\CE}{\mathcal{E}}
\newcommand{\CF}{\mathrm{CF}}
\newcommand{\CO}{\mathcal{CO}}
\newcommand{\CC}{\mathcal{C}}
\newcommand{\CL}{\mathcal{L}}
\newcommand{\bL}{\mathbb{L}}
\newcommand{\BL}{\mathbf{L}}
\newcommand{\A}{\mathbb{A}}
\newcommand{\CP}{\mathbb{C}P}
\newcommand{\bD}{\mathcal{D}}
\newcommand{\BA}{\mathbb{A}}
\newcommand{\bF}{\mathbb{F}}
\newcommand{\uF}{\underline{F}}
\newcommand{\bM}{\mathcal{M}}
\newcommand{\cF}{\mathscr{F}}
\newcommand{\cP}{\mathcal{P}}
\newcommand{\cS}{\mathcal{S}}
\newcommand{\cT}{\mathcal{T}}
\newcommand{\cX}{\mathcal{X}}
\newcommand{\cU}{\mathcal{U}}
\newcommand{\cH}{\mathcal{H}}
\newcommand{\Sh}{\mathfrak{Sh}}
\newcommand{\cY}{\mathcal{Y}}
\newcommand{\val}{\mathrm{val}}

\newcommand{\cb}{\mathfrak{b}}

\newcommand{\op}{\textrm{op}}

\newcommand{\SYZ}{\textrm{SYZ}}
\newcommand{\cpx}{\textrm{cpx}}
\newcommand{\Kah}{\textrm{K\"ah}}
\newcommand{\stab}{\textrm{stab}}

\newcommand{\Aut}{\mathrm{Aut}}
\newcommand{\aut}{\mathfrak{aut}}
\newcommand{\Hom}{\mathrm{Hom}}
\newcommand{\rank}{\mathrm{rank}}
\newcommand{\Dom}{\mathrm{Dom}}
\newcommand{\End}{\mathrm{End}}
\newcommand{\Fuk}{\mathrm{Fuk}}
\newcommand{\Ker}{\mathrm{Ker}\,}
\renewcommand{\Im}{\mathrm{Im}\,}

\newcommand{\GL}{\mathrm{GL}}
\newcommand{\SL}{\mathrm{SL}}
\newcommand{\SU}{\mathrm{SU}}
\newcommand{\Sp}{\mathrm{Sp}}

\newcommand{\cE}{\mathcal{E}}

\newcommand{\Jac}{\mathrm{Jac}}
\newcommand{\Span}{\mathrm{Span}}

\newcommand{\conn}{\nabla}
\newcommand{\Lieder}{\mathscr{L}}
\newcommand{\Laplacian}{\Delta}
\newcommand{\der}{\mathrm{d}}
\newcommand{\dd}[1]{\frac{\partial}{\partial #1}}
\newcommand{\dbar}{\bar{\partial}}

\newcommand{\contract}[2]{\iota_{#1} #2}
\newcommand{\dotprod}[2]{\left\langle #1 \, , \, #2 \right\rangle}
\newcommand{\pairing}[2]{\left( #1\, , \, #2 \right)}
\newcommand{\conj}[1]{\overline{#1}}
\newcommand{\WT}[1]{\widetilde{#1}}
\newcommand{\OL}[1]{\overline{#1}}
\newcommand{\oi}{\overline{\iota}}
\newcommand{\tr}{\mathrm{tr} \,}
\newcommand{\unu}{\underline{\nu}}
\newcommand{\Mor}{\mathrm{Mor}}

\newcommand{\nat}{\mathbb{N}}
\newcommand{\integer}{\mathbb{Z}}
\newcommand{\rat}{\mathbb{Q}}
\newcommand{\real}{\mathbb{R}}
\newcommand{\quat}{\mathbb{H}}
\newcommand{\oct}{\mathbb{O}}
\newcommand{\one}{\mathbf{1}}

\newcommand{\consti}{\mathbf{i}\,}
\newcommand{\conste}{\mathbf{e}}
\newcommand{\ev}{\mathrm{ev}}

\newcommand{\sphere}[1]{\mathbf{S}^{#1}}
\newcommand{\proj}{\mathbb{P}}
\newcommand{\bP}{\mathbb{P}}
\newcommand{\torus}{\mathbf{T}}
\newcommand{\moduli}{\mathcal{M}}
\newcommand{\cM}{\mathcal{M}}
\newcommand{\cA}{\mathcal{A}}
\newcommand{\cB}{\mathcal{B}}
\newcommand{\cC}{\mathcal{C}}
\newcommand{\cD}{\mathcal{D}}
\newcommand{\cO}{\mathcal{O}}
\newcommand{\cI}{\mathcal{I}}
\newcommand{\cJ}{\mathcal{J}}
\newcommand{\cL}{\mathcal{L}}
\newcommand{\pt}{\mathrm{pt}}
\newcommand{\AI}{A_\infty}
\newcommand{\MF}{\mathrm{MF}}
\newcommand{\Tw}{\mathrm{Tw}}
\newcommand{\LM}{\mathcal{LM}}
\newcommand{\QH}{\mathrm{QH}}
\newcommand{\CM}{\mathcal{LF}}
\newcommand{\mirror}{\mathrm{mirror}}
\newcommand{\barL}{\bar{L}}
\newcommand{\new}{\mathrm{new}}
\newcommand{\DT}{\mathrm{DT}}
\newcommand{\Coh}{\mathrm{Coh}}
\newcommand{\cd}{\check{\partial}}
\newcommand{\Id}{\mathrm{Id}}
\newcommand{\dt}[1]{\doubletilde{$#1$}}
\newcommand{\st}[1]{\tilde{#1}}
\newcommand{\Ham}{\mathrm{Ham}}
\newcommand{\Ad}{\mathrm{Ad}}

\title{Moment Lagrangians, unobstructedness and symplectic groupoids}

\author[Li]{Yan-Lung Leon Li}
\address{Center of Geometry and Physics, Institute for Basic Science (IBS), Pohang 37673, Korea}
\email{ylli@ibs.re.kr}


\begin{document}
\maketitle

\begin{abstract}
Moment Lagrangian $L_\mu$ is a Lagrangian in $T^*G^- \times Y^- \times Y$ associated to a Hamiltonian $G$-space $Y$ with a moment map $\mu$. In this paper, we prove that $L_\mu$ is tautologically unobstructed under mild assumptions on $Y$. 

As a key ingredient in the proof, we constructed a new symplectic groupoid structure on $T^*G^- \times Y^- \times Y$ over $G\times Y$ for which $L_\mu$ is simultaneously the unit and the fixed locus of the inversion,  which might be of independent interest. 
\end{abstract}

\section{Introduction}
Symplectic manifolds are the fundamental geometric objects underlying Hamiltonian mechanics, whose origin can be traced back to the classical works of Lagrange. Through the important works of Kostant, Kirillov and Souriau in the 1960s, it became clear that the natural geometric framework for Hamiltonian mechanics with symmetry group $G$ \footnote{Throughout the paper, $G$ is a compact, connected Lie group.} is provided by Hamiltonian $G$-spaces $(Y, \omega, \mu)$, where $(Y,\omega)$ is a symplectic manifold  with a Hamiltonian $G$-action in the sense that it admits a $G$-equivariant moment map $\mu$ satisfying the Hamilton's equation\footnote{See \ref{hamilton} for the precise formula.}.

Meanwhile, in \cite{Weinsteinsympgeom}, Weinstein proposed his famous symplectic creed, namely ``Everything is a Lagrangian submanifold", which led to his formulation of the symplectic ``category" $\mathrm{Symp}$. 
Roughly speaking, its objects are symplectic manifolds, its morphisms are Lagrangian correspondences between them, and the composition of morphisms is given by geometric composition of correspondences. In this framework, Hamiltonian $G$-spaces $(Y, \omega, \mu)$ can be incorporated naturally through their moment Lagrangians
$$L_\mu \coloneqq \{(g, \mu(y), y, g\cdot y) | g\in G; y \in Y\} \subseteq (T^*G)^- \times Y^- \times Y \footnote{In the introduction, we write $T^*G$ for simplicity and easier comparison with existing literature; in later sections, we will replace $T^*G$ by $G \times \mathfrak{g}^*$ 
to avoid confusions, since its symplectic form $\Omega$ will be chosen to be correspond to $-\omega_{can}$; see Definition \ref{def:momentlag} for details.}$$
as a Lagrangian correspondence from $T^*G$ to  $Y^- \times Y$, providing a realization of the action as a morphism in $\mathrm{Symp}$. This perspective has led to profound applications in quantization and semi-classical analysis; see, for instance, \cite{GSsemiclassical} and the references therein.

On the other hand, in his ICM address \cite{Teleman}, Teleman proposed using $L_\mu$ to construct a 
topological $G$-action on the Fukaya category $\mathrm{Fuk}(Y)$ of $Y$. A key intermediate step is to use $L_\mu$ 
to induce a functor from the wrapped Fukaya category of $T^*G$ to the Fukaya category of $Y^-\times Y$, 
which the latter acts on $\mathrm{Fuk}(Y)$ via geometric transformations. Teleman's proposal has since inspired substantial interests in the Lagrangian Floer theory of $L_\mu$ and its applications, resulting in important developments in the monotone setting, including \cite{EL, BCL, Chowmomentcorr} to name a few.

Meanwhile, to the best of the author's knowledge, the Floer theory of $L_\mu$ for more general symplectic manifolds, in terms of foundational results of Fukaya, Oh, Ohta and Ono in their series of works on Lagrangian Floer theory (e.g. \cite{FOOO,foookuranishi}), seems to remain largely unexplored. In particular, a fundamental question is to determine whether or not $L_\mu$ is (Floer-theoretically) unobstructed, in the sense that its (Lagrangian) Floer complex $CF(L_\mu)$ is an unobstructed $A_\infty$ algebra (in particular a chain complex) after some bulk or boundary deformations\footnote{See Remark \ref{bdydeform} for the notion of boundary deformation by bounding cochains.} so that its Floer cohomology $HF(L_\mu)$ is well-defined after such deformations. Establishing unobstructedness is therefore a necessary first step towards any applications of the Floer theory of $L_\mu$.

In this work, we prove the tautogical unobstructedness of $L_\mu$ under some mild assumptions as follows:

\begin{theorem}\label{momentlagrangianunobstructed}
Given a closed, connected, spin Hamiltonian $G$-space $(Y, \omega, \mu)$, its moment Lagrangian $L_\mu$ is tautologically unobstructed. 
\end{theorem}

\begin{remark}
In \cite[Lemma 6.1.1]{EL}, Evans and Lekili established the monotonicity of $L_\mu$ under the additional assumptions that $Y$ is simply-connected and monotone, in order to construct an $A_\infty$ functor from the wrapped Fukaya category of $T^*G$ to the monotone Fukaya category $\mathrm{Fuk}(Y^- \times Y)$. This has played a pivotal role in computing $\mathrm{Fuk}(Y)$ in terms of determining its split generators. 

From this perspective, Theorem \ref{momentlagrangianunobstructed} shows that the unobstructedness of $L_\mu$ persists beyond monotone setting. Moreover, combined with Fukaya's result in \cite[Theorem 1.7]{Fukaya-corr}, it yields an $A_\infty$ functor between their (compact, unobstructed) Fukaya categories. It would be interesting to see if this functor would also be useful in computing $\mathrm{Fuk}(Y)$ in the sense of \cite{AFOOO}.
\end{remark}

\begin{remark}
    This result should not be confused with the unobstructedness result of Xiao in \cite{xiao2}. In the loc. cit, the term ``moment Lagrangian correspondence" refers to the correspondence between a Hamiltonian $G$-space $Y$ and its symplectic quotient $Y \sslash_0 G$ (which is called ``moment level correspondence" in \cite{lauleungli}).
\end{remark}

As an application of Theorem \ref{momentlagrangianunobstructed}, together with Fukaya's unobstructedness theorem on geometric transformations, we obtain the unobstructedness of the following Lagrangians:

\begin{corollary}\label{cor: exunoblag}
    \begin{enumerate}
        \item Given a closed, connected, spin Hamiltonian $G$-space $(Y, \omega, \mu)$ such that $G$ acts freely on $\mu^{-1}(0)$, the following Lagrangian is unobstructed: $$\Delta_\mu \coloneqq  \{(y, g\cdot y) \in Y^- \times Y | g\in G; y\in \mu^{-1}(0)\} \subset Y^- \times Y.$$
        \item Given a coadjoint orbit $\mathcal{O}_\lambda$ of $G$, the following Lagrangian is unobstructed: $$L_{\mathcal{O}_\lambda} \coloneqq  \{(g, \Ad^*_g(\xi)) \in G \times \mathfrak{g}^* | g\in G; \xi\in \mathcal{O}_\lambda\}\subset G \times \mathfrak{g}^*.$$
    \end{enumerate}
\end{corollary}

\begin{remark}
In the monotone setting, $\Delta_\mu$ was also studied by Evans and Lekili in \cite[Section 6.2]{EL}, where $\Delta_\mu$ is assumed to be monotone. 
%because monotonicity of Lagrangians needs not be preserved under geometric transformations
From this perspective, (1) can be viewed as showing that $\Delta_\mu$ is in fact unobstructed in general. 
%which follows from Theorem \ref{unobstructedness} that unobstructedness is preserved under geometric transformations.
\end{remark}

Returning to the proof of Theorem \ref{momentlagrangianunobstructed}, the key idea is to realize $L_\mu$ as the fixed locus of an anti-symplectic involution $i$. Its unobstructedness then follows from a theorem of Fukaya, Oh, Ohta and Ono \cite[Corollary 1.7]{FOOOinvol}. 
Interestingly, $L_\mu = \mathrm{Fix}(i)$ admits a broader interpretation in terms of symplectic groupoids, which we briefly recalled below.

Symplectic groupoid was introduced by Weinstein in \cite{weinsteinsympgroupoid}, with a more detailed treatment appearing in \cite{sympgroupoid}. Roughly speaking, a symplectic groupoid is a Lie groupoid $X$ over $M$ such that $X$ admits a multiplicative symplectic structure integrating the Poisson structure on $M$\footnote{See Definition \ref{def:sympgroupoid} for the precise definition.}.  The notion arises naturally in the study of symplectic realizations and quantization of Poisson manifolds. 

In this work, we construct a new example of symplectic groupoid on $X \coloneqq (T^*G)^- \times Y^- \times Y$ for which the moment Lagrangian $L_\mu$ is the unit, a result which might be of independent interest, as follows:

\begin{theorem}(= Theorem \ref{twisted-groupoid}) \label{intro:twisted-groupoid}
    Given a Hamiltonian $G$-space $(Y, \omega, \mu)$, there exists a symplectic groupoid structure on $X$ over $G \times Y$
 such that $L_\mu$ is the unit and also the fixed locus of the anti-symplectic involution $i$:
    $$L_\mu = \underline{e} = Fix(i).$$
\end{theorem}

\begin{remark}
    The equality $L_\mu = \underline{e}$ already appears in \cite{Weinsteinsympgeom}. The theorem above provides a new interpretation of this from the perspective of symplectic groupoids.
\end{remark}

\begin{remark}
As another consequence of \cite[Corollary 1.6]{FOOOinvol} applied to $L_\mu = Fix(\underline{e})$, $HF(L_\mu)$ admits a graded-commutative product induced by $m_2$ (after suitably twisted by some signs). 

Another interesting way to see this is via $L_\mu  = \underline{e}$ in Theorem \ref{intro:twisted-groupoid}: graded commutativity of $HF(L_\mu)$ follows from a 
general Eckmann-Hilton-type argument applied to $\End(\underline{e})$, where $\underline{e}$ is a unit object in the monoidal Fukaya category of a symplectic groupoid, due to Pascaleff in \cite[Section 5.1]{Pascaleffmonoidal}. This provides another instance in which structural properties of symplectic groupoids implies graded commutativity of Floer cohomology of units.
\end{remark}

Returning to the proof of Theorem \ref{intro:twisted-groupoid}, the groupoid structure is given explicitly in 
Theorem \ref{twisted-groupoid}, which can be checked to satisfy the axioms of Lie groupoid \footnote{See Appendix \ref{appendix:liegroupoid} for the axioms.}. The multiplicativity of the symplectic form follows from the properties of moment map, 
although the derivation is somewhat intricate, as demonstrated in the proof of Theorem \ref{twisted-groupoid}. Perhaps the most mysterious part of the construction is the groupoid structure itself, even to the author.

\subsection*{Organization of this paper} In Section 2, we recall the notions of moment Lagrangian and symplectic groupoid. Then we prove Theorem \ref{intro:twisted-groupoid} (= Theorem \ref{twisted-groupoid}) on the construction of the symplectic groupoid on $X$.

In Section 3, we first briefly recall the notions of Lagrangian Floer complexes and their unobstructedness. Then we prove Theorem \ref{momentlagrangianunobstructed} on the unobstructedness of $L_\mu$. 
Afterwards, we prove Corollary \ref{cor: exunoblag} on unobstructedness of some geometric transformations by $L_\mu$.

In Appendix \ref{appendix:symplecticform}, we recall an explicit formula for the symplectic form $\Omega$ on $G \times \mathfrak{g}^*$.

In Appendix \ref{appendix:liegroupoid}, we recall the notions of Lie groupoid to clarify the choices and conventions for this paper. Other conventions are listed in the following subsection.

\subsection*{Notations and Conventions}
Throughout the whole paper, we will generally follow the conventions in \cite{CFM}. Some of the relevant ones are as follows:
\begin{enumerate}
\item For any smooth manifold $M$, its cotangent bundle $T^*M$ has a canonical symplectic form $\omega_{can} = -d\lambda_{can}$, i.e. locally of the form $\sum_{i} dq_i \wedge dp_i$ with respect to local coordinates $(q_i)$ of $M$ and cotangent fiber coordinates $(p_i)$.

%\item Given a (real) Lie group $G$, denote (the dual space of) its Lie algebra by $\mathfrak{g}$ (resp. $\mathfrak{g}$*). We will always identify its cotangent bundle $T^*G$ with $G \times \mathfrak{g}^*$ via left translation, i.e. 
%$$G \times \mathfrak{g}^* \xrightarrow{l} T^*G$$
%$$(g, \xi) \mapsto \mathcal{L}^*_g(\xi)\in T^*_gG$$
%where $\mathcal{L}^*_g: \mathfrak{g}^* \rightarrow T^*_gG$ is the dual map of the left translation $\mathcal{L}^{-1}_{g}: T_gG\rightarrow \mathfrak{g}$.

\item Given a Hamiltonian $G$-space $(Y, \omega, \mu)$, the moment map $\mu$ is a $G$-equivariant map $\mu: Y \rightarrow \mathfrak{g}^*$, where $\mathfrak{g}^*$ is the dual of the Lie algebra of $G$ endowed with coadjoint $G$-action, satisfying the following Hamilton's equation:
\begin{equation}\label{hamilton}
  \iota_{v^\#}\omega = \langle d\mu(-), v \rangle, \forall v\in \mathfrak{g}.
\end{equation}

where $v^\#$ is the infinitesimal generator of the $G$-action on $Y$ associated to $v\in \mathfrak{g}$ defined as $v^\#(y) = \frac{d}{dt}|_{t=0} \exp(-tv)\cdot y$ for any $y\in Y$, which gives rise to a Lie algebra homomorphism from $\mathfrak{g}$ to the Lie algebra of vector fields on $Y$.

Under this convention, the $G$-equivariance of $\mu$ is equivalent to that the comoment map $\mu^*: \mathfrak{g} \rightarrow C^\infty(Y)$, given by $\mu^*(v) \coloneqq \langle \mu(-), v\rangle$, is a Lie algebra homomorphism.


%\item Given a smooth $G$-space $M$, i.e. a compact connected Lie group $G$ which acts smoothly (on the left) on a smooth manifold $M$, it lifts to a Hamiltonian action on $(T^*M, \omega_{can})$ with a (canonical) moment map $\mu$, which under the above convention is defined as 
%$\mu(\xi)(v) = \langle\xi, v^\#\rangle.$

%In particular, when $M = G$ acted by left multiplication, the moment map $\mu$ on $T^*G \cong G \times \mathfrak{g}^*$ is given by $\mu(g, \xi)\mapsto \Ad^*_g(\xi)$.

\end{enumerate}

\section*{Acknowledgments}
The author would like to thanks Hanwool Bae, 	
Sam Bardwell-Evans, Siu-Cheong Lau, Conan Leung, Yong-Geun Oh and Yutung Yau for helpful discussions on various stages of this project. He would also like to thank Ziming Ma for the hospitality in the International Conference on Mathematics Inspired by String Theory, in which part of this work was done and presented.
This work was supported by the Institute for Basic Science (IBS-R003-D1).

\section{Moment Lagrangians and symplectic groupoids}


In this section, we first recall the notion of moment Lagrangian and symplectic groupoid, and then construct a new example of symplectic groupoid whose unit is the moment Lagrangian.

The precise definition of moment Lagrangian is as follows:

\begin{defn}\label{def:momentlag}
    Given a Hamiltonian $G$-space $(Y, \omega)$ with moment map $\mu$, its moment Lagrangian $L_\mu$ is defined as 
    $$L_\mu = \{(g, \mu(y), y, g\cdot y) | g\in G; y \in Y\} \subseteq ((G \times \mathfrak{g}^*)^- \times Y^- \times Y, (-\Omega) \oplus (-\omega) \oplus \omega).$$
    where $Y^-$ is the symplectic manifold $(Y, -\omega)$, and $\Omega$ is the symplectic form on $G \times \mathfrak{g}^*$ corresponding to $-\omega_{can}$ under trivialization $G \times \mathfrak{g}^* \cong T^*G$ by left translation. \footnote{See Appendix \ref{appendix:symplecticform} for further details, including an explicit formula of $\Omega$.}
\end{defn}

\begin{remark}
This choice of symplectic structure $\Omega$ on $G \times \mathfrak{g}^*$ is for the compatibility with the chosen conventions for symplectic groupoids. See Appendix \ref{appendix:symplecticform} for further details.
\end{remark}

%\begin{example} Consider $X = (T^*G, \omega_{can})$ with lifted Hamiltonian $G$-action by left multiplication and moment map $\mu = \Ad^*$, its moment Lagrangian is 
%$$L_\mu = \{((g, \Ad^*_h(\xi)), (h, \xi), (gh, \xi) \in (T^*G)^- \times (T^*G)^- \times (T^*G)| g \in G, (h, \xi)\in T^*G\}.$$
    
%\end{example}

We now recall the notion of symplectic groupoids as follows:

\begin{defn}\label{def:sympgroupoid}
A symplectic groupoid is a Lie groupoid
%\footnote{\label{a}See Definition \ref{liegroupoid}, Example \ref{pairgroupoid} and \ref{actiongroupoid}  for their precise definitions.} 
$(X \xRightarrow[t]{s} M, m, e, i)$, where $(X, \omega)$ is a symplectic manifold such that $\omega$ is multiplicative: 
\[
m^*\omega = pr_1^*\omega+ pr_2^*\omega \in \Omega^2(X^{(2)}),
\]
where $pr_j: X^{(2)} \rightarrow X$ be the projection to $j$-th component.
\end{defn}

%\begin{remark}\label{poisson}
%A symplectic groupoid structure $(X, \omega) \xRightarrow[t]{s} M$ induces a unique Poisson structure on $M$ such that $s$ is an anti-Poisson map and $t: X \rightarrow M$ is a Poisson map under our conventions. 
%See e.g. \cite[Theorem 14.10]{CFM} for more details.
%\end{remark}

\begin{prop} \label{lagrangiansubmanifold}
Given a symplectic groupoid $(X, \omega) \xRightarrow[t]{s} M$, the following are Lagrangians:
\begin{enumerate}
\item $\underline{m} := Graph(m) = \{(p,q,p\cdot q) \ |\,  (p,q)\in X^{(2)}\} \subseteq X^- \times X^- \times X,$
\item $\underline{e} := e(M)=\{e(x) |\,  x\in M\} \subseteq X,$
\item $\underline{i} := Graph(i) =\{(p,p^{-1})) \in X \times X|\,  p\in X \}\subseteq X \times X.$ 
\end{enumerate}
\end{prop} 

\begin{proof}
See e.g. \cite[Prop. 14.9]{CFM}.
\end{proof}

We now state the main result of this paper, which constructs a new example of symplectic groupoid whose unit is the moment Lagrangian:

\begin{theorem}\label{twisted-groupoid}
    Given a Hamiltonian $G$-space $(Y, \omega, \mu)$, there exists a symplectic groupoid structure on $(X \coloneqq (G \times \mathfrak{g}^*)^- \times Y^- \times Y, \omega_X \coloneqq ((-\Omega)\oplus (-\omega) \oplus \omega))$:
    $$((X , \omega_X) \xRightarrow[t]{s} G \times Y, m, e, i),$$
given by the following:
    \begin{itemize}
        \item  $s(g, \xi, y, z) = (g, g^{-1}\cdot z)$; $t(g, \xi, y, z) = (g, y)$;
        \item $(g_1, \xi_1, y_1, z_1) \cdot (g_2, \xi_2, y_2, z_2) = (g, \xi_1 + \xi_2 - \mu(y_2), y_1, z_2)$, where 
        $$g_1 = g_2 = g; z_1 = g\cdot y_2;$$
        \item $e(g, x) = (g, \mu(x), x, g\cdot x)$;
        \item $i(g, \xi, y, z) = (g, \mu(g^{-1}\cdot z)+ \mu(y) - \xi, g^{-1}\cdot z, g\cdot y)$.
    \end{itemize}

     In particular, $L_\mu$ is the unit and also the fixed locus of the anti-symplectic involution $i$:
    $$L_\mu = \underline{e} = Fix(i).$$
\end{theorem}



\begin{remark}
Combining Theorem \ref{twisted-groupoid} with Proposition \ref{lagrangiansubmanifold}, we obtain an alternative proof of the fact that the moment Lagrangian $L_\mu$ is a Lagrangian of $(X, \omega_X)$.
\end{remark}

\begin{proof}
It is straightforward to check that the above maps satisfy the axioms of a Lie groupoid, and that $L_\mu = \underline{e} = Fix(i)$. It remains to check that $\omega_X$ is multiplicative, which will be done by definition as follows.

    First of all, the tangent space of $X^{(2)}$ at $(x_1, x_2) \coloneqq ((g, \xi_1, y_1, z_1), (g, \xi_2, y_2, z_2))$ is given by 
    $$T_{(x_1, x_2)}X^{(2)} = \{((Z, \eta_1, v_1, v'_1), (Z, \eta_2, v_2, v'_2))| Z\in \mathfrak{g}; v'_1 = g_* v_2 + (y_2)_* Z\}.$$
    where $g_*: T_{y}Y \rightarrow T_{g\cdot y}Y$ is the differential of the action by $g$ and $y_*: \mathfrak{g} \rightarrow T_{y}Y$ is the infinitesimal generator $Z^\#$ evaluated at $y$.

    Also, the differential of $m$ at $(x_1, x_2)$ is given by
 $$dm_{(x_1, x_2)}((Z, \eta_1, v_1, v'_1), (Z, \eta_2, v_2, v'_2))=(Z, \eta_1 + \eta_2 - d\mu(v_2), v_1, v'_2).$$

    Therefore, the symplectic form is multiplicative if and only if for any 
    $$((Z, \eta_1, v_1, v'_1), (Z, \eta_2, v_2, v'_2)), (W, \xi_1, w_1, w'_1), (W, \xi_2, w_2, w'_2))\in T_{(x_1, x_2)}X^{(2)},$$
    the following holds:
\begin{align*}
    &\omega_X((Z, \eta_1 + \eta_2 - d\mu(v_2), v_1, v'_2), (W, \xi_1 + \xi_2 - d\mu(w_2), w_1, w'_2))\\
    = &\omega_X((Z, \eta_1, v_1, v'_1), (W, \xi_1, w_1, w'_1)) + \omega_X((Z, \eta_2, v_2, v'_2), (W, \xi_2, w_2, w'_2)).
\end{align*}
Using \ref{symplecticform}, the above is equivalent to the following:
\begin{align*}
&-\langle d\mu(v_2), \mathcal{L}^{-1}_{g}(W) \rangle + \langle d\mu(w_2), \mathcal{L}^{-1}_{g}(Z) \rangle + \langle \mu(y_2), [\mathcal{L}^{-1}_{g}(Z), \mathcal{L}^{-1}_{g}(W)] \rangle \\
= &-\omega(g_*(v_2), (y_2)_*W) - \omega((y_2)_*Z, g_*(w_2)) - \omega((y_2)_*Z, (y_2)_*W).
\end{align*}
The above equality holds if the following identities hold:
\begin{align*}
\langle d\mu(v), \mathcal{L}^{-1}_{g}(W) \rangle &= \omega(g_*(v), (y_2)_*W);\\
\langle \mu(y_2), [\mathcal{L}^{-1}_{g}(Z), \mathcal{L}^{-1}_{g}(W)] \rangle &= - \omega((y_2)_*Z, (y_2)_*W).
\end{align*}
The first identity follows from the Hamilton's equation \ref{hamilton}, and the second identity follows from that the comoment map $\mu^*$ is a Lie algebra homomorphism.
\end{proof}

\section{Applications to unobstructedness}

In this section, we provide some applications of Theorem \ref{twisted-groupoid} to obtain results on (Lagrangian) Floer theory of $L_\mu$ and some geometric transformations by it. 

\subsection{Floer theory and unobstructedness}
In this section, we briefly review the notions of Floer complex and its unobstructedness. Readers are referred to e.g. \cite{Fukaya-corr, FOOO,FOOOinvol} for more details.

Given a closed (or tame) symplectic manifold $(X, \omega)$ and a closed, connected, oriented and relatively spin Lagrangian $L$ in $X$, the Floer complex $CF(L)$ of $L$ is a filtered $A_\infty$ algebra $(CF(L), \{m_k\}_{k\geq 0})$ over $\Lambda_0$
\footnote{where $\Lambda_0 = \{\sum_{i=0}^\infty a_i T^{\lambda_i}| a_i \in \mathbb{R}; 0=\lambda_0<\lambda_1<\cdots; \lambda_i \rightarrow \infty\}$ is the (universal) Novikov ring defined in \cite[Definition 2.1]{Fukaya-corr}.}, whose structure maps $m_k: CF(L)^{\otimes k} \rightarrow CF(L)$ are defined by counting pseudo-holomorphic discs with boundary on $L$ and satisfy the $A_\infty$ relations. In particular, the following equation holds:
$$m_1 \circ m_1 = m_0(1) \cdot \Id: CF(L) \rightarrow CF(L).$$
Therefore, if the curvature term $m_0(1)$ vanishes, then $(CF(L), m_1)$ is a chain complex with cohomology $HF(L)$ being the Floer cohomology of $L$.

In general, $m_0(1)$ may not vanish, which is an obstruction to $(CF(L), m_1)$ being a chain complex. However, $CF(L)$ could become unobstructed after deforming the $A_\infty$ algebra by a bounding cochain $b$ defined as follows:

\begin{defn}\label{unobstructedlag}
$L$ is called unobstructed if $(CF(L), \{m_k\})$ is unobstructed as a filtered $A_\infty$ algebra, i.e. there exists a bounding cochain $b\in CF^1(L; \Lambda_+)$\footnote{where $\Lambda_+ = \{\sum_{i=1}^\infty a_i T^{\lambda_i}| a_i \in \mathbb{R}; \lambda_i > 0\}$ is the (unique) maximal ideal of $\Lambda_0$.} such that the Maurer-Cartan equation holds:
$$\sum_{k=0}^{\infty} m_k(b, \ldots, b) = 0.$$  

In particular, $L$ is called tautologically unobstructed if $0$ is a bounding cochain, i.e. $m_0(1) = 0$.
\end{defn}

\begin{remark}\label{bdydeform}
    Any bounding cochain $b$ defines a deformed $A_\infty$ algebra $(CF(L), \{m_k^b\})$ with $m_0^b(1) \coloneqq \sum_{k=0}^{\infty} m_k(b, \ldots, b)  = 0$ by assumption, and hence $(CF(L), m_1^b)$ is a chain complex with cohomology $HF(L, b)$ being the deformed Floer cohomology of $(L, b)$. See \cite[Definition 6.3]{FOOOinvol} and the reference therein for more details.
\end{remark}

\subsection{Unobstructedness of moment Lagrangians}

In this section, we prove Theorem \ref{momentlagrangianunobstructed} on tautological unobstructedness of the moment Lagrangian using Theorem \ref{twisted-groupoid} as follows:

\begin{proof}[Proof of Theorem \ref{momentlagrangianunobstructed}]
    The assumptions guarantee that $X$ is tame and spin; also, the moment 
    Lagrangian $L_\mu \cong G \times Y$ is a closed, connected, oriented and relatively spin Lagrangian of $(X, \omega_X)$. By Theorem \ref{twisted-groupoid}, $L_\mu = Fix(i)$ is the fixed locus of an anti-symplectic involution $i$; also, the induced map $\iota_*: \pi_1(L_\mu)\rightarrow \pi_1(X)$ of the inclusion $\iota: L_\mu \hookrightarrow X$ is injective, since $\iota$ admits a left inverse given by 
    $e \circ t: X \rightarrow L_\mu$. Therefore, by \cite[Corollary 1.6]{FOOOinvol} (in which $b$ can be chosen to be zero as shown in its proof), $L_\mu$ is tautologically unobstructed.
\end{proof}




\subsection{Unobstructedness of geometric transformations}

In this section, we provide more examples of unobstructed Lagrangians as geometric transformations of $L_\mu$ as an unobstructed Lagrangian correspondence.

Firstly, we recall the notions of Lagrangian correspondences and geometric transformations by them, which are first introduced by Weinstein in his symplectic category formalism in \cite{sympcategory}:

\begin{defn}
Given a Lagrangian correspondence $\mathcal{L}$ from $(M, \omega_M)$ to $(N, \omega_N)$, i.e. a Lagrangian in $(M^- \times N, -\omega_M\oplus\omega_N)$. For any Lagrangian $L \subseteq (M, \omega_M)$, the geometric transformation of $L$ by $\mathcal{L}$, denoted as $\mathcal{L} \circ L$, is a subset of $N$ defined as
\begin{equation}
\mathcal{L} \circ L = pr_{N}((L \times \mathcal{L})\cap (\Delta_{M} \times N)),
\end{equation}
where $pr_{N} :M \times M \times N \rightarrow N$ is the natural projection.\\

$\mathcal{L} \circ L$ is called embedded if 
   $(L \times \mathcal{L})\cap (\Delta_{M} \times N)$ is a transverse intersection in $M \times M \times N$, and    
$pr_{N}$ restricted to a smooth embedding on this intersection. If so, then $\mathcal{L} \circ L$ is an embedded Lagrangian of $(N, \omega_N)$.
\end{defn}

Examples of geometric transformations by $L_\mu$ (and its transpose $L^t_\mu$) are given as follows:

\begin{example}\label{example:shift-diagonal}
The moment Lagrangian $L_\mu \subseteq (G \times \mathfrak{g}^*)^- \times Y^- \times Y$ is a Lagrangian correspondence from $(G \times \mathfrak{g}^*, \Omega)$ to $(Y^- \times Y, -\omega \oplus \omega)$. Let $L \coloneqq G \times {0} \subseteq (G \times \mathfrak{g}^*)$ be the zero section, then the geometric transformation of $L$ by $L_\mu$ is given by 
$$\Delta_\mu \coloneqq L_\mu \circ L =  \{(y, g\cdot y) \in Y^- \times Y | g\in G; y\in \mu^{-1}(0)\},$$
which is embedded if $G$ acts freely on $Z \coloneqq \mu^{-1}(0)$ with smooth symplectic quotient $\bar{Y} \coloneqq Y\sslash_0 G = Z/G$. In this case, $\Delta_\mu$ can also be written as a fiber product $Z \times_{\bar{Y}} Z$.
\end{example}    

\begin{example}\label{example:orbit-lagrangian}
    Let $(Y = \mathcal{O}_\lambda, \omega)$ be a coadjoint orbit of $G$ with the Kirillov-Kostant-Souriau symplectic form $\omega$, 
    for which $G$ acts in a Hamiltonian fashion via coadjoint action with moment map $\mu = \iota: \mathcal{O}_\lambda \hookrightarrow \mathfrak{g}^*$. 
    This time, we consider the transpose of the moment Lagrangian:
    $$L^t_\mu \coloneqq \{(y, g\cdot y, g, \mu(y))) \subseteq ((\mathcal{O}_\lambda^- \times \mathcal{O}_\lambda)^- \times (G \times \mathfrak{g}^*), -(-\omega\oplus\omega)\oplus\Omega)$$ 
    as a Lagrangian correspondence from $\mathcal{O}_\lambda^- \times \mathcal{O}_\lambda$ to $G \times \mathfrak{g}^*$. Let $L = \Delta_{\mathcal{O}_\lambda}$ be the diagonal Lagrangian, 
    then the geometric transformation of $L$ by $L^t_\mu$ is given by the orbit Lagrangian:
    $$L_{\mathcal{O}_\lambda} \coloneqq L^t_\mu \circ L =  \{(g, \Ad^*_g(\xi)) \in G \times \mathfrak{g}^* | g\in G; \xi\in \mathcal{O}_\lambda\},$$
    which is embedded since $\mu$ is injective.
\end{example}

We now recall the following result of Fukaya on preservation of unobstructedness of Lagrangians under geometric transformations below:

\begin{theorem}\cite[Theorem 1.5]{Fukaya-corr} \label{unobstructedness}
Given two closed or tame symplectic manifolds $(M, \omega_M)$,
$(N, \omega_N)$ and a closed, connected, oriented and relatively spin Lagrangian correspondence $\mathcal{L} \subseteq (M^-\times N, \omega_M\oplus \omega_N)$ which is unobstructed with a bounding cochain $b_{\mathcal{L}}$, then for any closed, connected, oriented and relatively spin unobstructed Lagrangian $L \subseteq (M, \omega_M)$ with a bounding cochain $b$ such that the geometric transformation
$L' \coloneqq \mathcal{L} \circ L$ is also closed, connected, oriented and relatively spin, then $L'$ is also unobstructed with a bounding cochain $b'$ canonically determined by $b$ and $b_{\mathcal{L}}$.
\end{theorem}

\begin{remark}
This can be generalised to weak bounding cochains, i.e. $b$ and $b_{\mathcal{L}}$ are weak bounding cochains, then $b'$ is also a weak bounding cochain satisfying an identity of their disc potentials. Interested readers are referred to \cite[Proposition 2.10]{lauleungli} for details. Meanwhile, Theorem \ref{unobstructedness} would suffice for applications in this paper.
\end{remark}

\begin{remark}
    Readers are cautioned that tautological unobstructedness needs not be preserved under geometric transformations: $b = 0 = b_{\mathcal{L}}$ does not imply $b' = 0$, unless the cyclic element in the correspondence tri-module is closed under the tri-module differential. See also \cite[Lemma 7.21]{gao} for relevant discussion in the exact case. 
\end{remark}

We are now ready to prove Corollary \ref{cor: exunoblag} on the unobstructedness of the geometric transformations 
in Examples \ref{example:shift-diagonal} 
and \ref{example:orbit-lagrangian} are unobstructed:

\begin{proof}[Proof of Corollary \ref{cor: exunoblag}]
For (1), the assumptions guarantee that $L_\mu$ is tautologically unobstructed by 
Corollary \ref{momentlagrangianunobstructed}; also, $Z$ is closed and connected 
(by Atiyah-Guillemin-Sternberg convexity theorem), so as $Z \times_{\bar{Y}} Z \subseteq Y^- \times Y$. 
Therefore, since $L$ is exact (and hence tautologically unobstructed), Theorem \ref{unobstructedness} implies that $\Delta_\mu$ is unobstructed.\\

For (2), since $L^t_\mu$ is also tautologically unobstructed and $\Delta_{\mathcal{O}_\lambda}$ is tautologically unobstructed by \cite[Theorem 1.9]{FOOOinvol}, Theorem \ref{unobstructedness} implies that $L_{\mathcal{O}_\lambda}$ is unobstructed.
\end{proof}



\appendix

\section{Symplectic structure of $G \times \mathfrak{g}^*$} \label{appendix:symplecticform}
Here we provide explicit formula for the symplectic form $\Omega$ on $G \times \mathfrak{g}^*$ under the identification with $(T^*G, -\omega_{can})$ via left translation, following \cite[Example 14.7]{CFM}. This was used in the proof of Theorem \ref{twisted-groupoid}. 

Given any compact connected Lie group $G$, there exists a diffeomorphism 
\begin{align*}
l: G \times \mathfrak{g}^* &\rightarrow T^*G\\
(g, \xi) &\mapsto \mathcal{L}^*_g(\xi),
\end{align*}
where $\mathcal{L}^*_g: \mathfrak{g}^* \rightarrow T^*_gG$ is dual to the inverse of the left translation $\mathcal{L}^{-1}_{g}: T_gG\rightarrow \mathfrak{g}$.

The symplectic form $\Omega$ on $G \times \mathfrak{g}^*$ is given by $\Omega = -l^*(\omega_{can})$. More explicitly, for any $(g, \xi) \in G \times \mathfrak{g}^*$ and $(Z, \eta), (W, \zeta) \in T_{(g, \xi)}(G \times \mathfrak{g}^*) \cong T_gG \times \mathfrak{g}^*$, we have
\begin{align}\label{symplecticform}
\Omega((Z, \eta), (W, \zeta)) = \langle \eta, \mathcal{L}^{-1}_{g}(W) \rangle - \langle \zeta, \mathcal{L}^{-1}_{g}(Z) \rangle - \langle \xi, [\mathcal{L}^{-1}_{g}(Z), \mathcal{L}^{-1}_{g}(W)] \rangle.
\end{align}

Using this symplectic form $\Omega$, the action groupoid $G \times \mathfrak{g}^*$ of the coadjoint $G$-action on $\mathfrak{g}^*$ is a symplectic groupoid $(G \times \mathfrak{g}^*, \Omega) \xRightarrow[Ad^*]{pr} \mathfrak{g}^*$. See e.g. \cite[Example 14.7]{CFM} for more details.

\section{Lie groupoids}\label{appendix:liegroupoid}
%We refer the readers to [Lectures on Poisson Geometry] for details.
We follow the convention of Lie groupoids in \cite{CFM}, recalled below for reader's convenience:
\begin{defn} \cite[Definition 13.1]{CFM} \label{liegroupoid}
A Lie groupoid consists of a pair of smooth manifolds $(X, M)$ and a five-tuple of smooth maps $(s,t, m, e, i)$, where 
\begin{itemize}
\item (source and target maps) $s,t: X \rightarrow M$ are surjective submersions.
\item (multiplication map) $m: X^{(2)} \rightarrow X$ is a smooth map, where 
$$ X^{(2)} \coloneqq \{(p,q) \in X \times X: s(p) = t(q)\},$$
which is smooth since $s$ and $t$ are submersions. Denote $p\cdot q \coloneqq m(p,q)$.
\item (unit map) $e: M \rightarrow X$ is a smooth embedding.
\item (inversion map) $i: X \rightarrow X$ is a smooth involution, i.e. $i\circ i=\Id$. We write $p^{-1} \coloneqq i(p)$.
\end{itemize}
such that for each $(p,q),(q,r) \in X^{(2)}$ and $x\in M$, the following conditions are satisfied:
\begin{itemize}
\item $e$ is a section of $s$ and $t$: $s(e(x)) = x = t(e(x))$.
\item $e(t(p)) \cdot p  = p =  p \cdot e(s(p)) $
\item $s(p \cdot q) = s(q) ;  t(p \cdot q) =t(p)$
\item $(p \cdot q) \cdot r = p \cdot (q \cdot r)$.
\item $s(p^{-1}) = t(p) ; t(p^{-1}) = s(p)$.
\item $p \cdot p^{-1} = e(t(p)); p^{-1} \cdot p = e(s(p))$.
 
\end{itemize}
For simplicity, we will denote a Lie groupoid by $(X \xRightarrow[t]{s} M, m, e, i)$.\\
\end{defn}


\bibliographystyle{plain}
\bibliography{geometry}
	
\end{document}
